\documentclass[11pt]{article}

\usepackage[margin=0.65in]{geometry}
\usepackage{amsmath,amssymb,mathtools}
\usepackage{bm}
\usepackage{booktabs}
\usepackage[hidelinks]{hyperref}
\usepackage{enumitem}

\setlist[itemize]{leftmargin=1.4em,itemsep=0.2em,topsep=0.2em}
\setlist[enumerate]{leftmargin=1.6em,itemsep=0.2em,topsep=0.2em}
\setlength{\parindent}{0pt}
\setlength{\parskip}{0.55em}

\newcommand{\data}{d}
\newcommand{\signal}{s}
\newcommand{\synth}{s'}
\newcommand{\target}{\tilde{s}}
\newcommand{\net}{f_\theta}
\newcommand{\mean}{\mu}
\newcommand{\logvar}{\ell}
\newcommand{\std}{\sigma}
\newcommand{\pix}{r}
\newcommand{\chan}{c}
\newcommand{\grid}{\Omega}
\newcommand{\batch}{\mathcal{B}}
\newcommand{\E}{\mathbb{E}}
\newcommand{\R}{\mathbb{R}}

\title{Moment Network Training: Legacy Std Calculation and Current Variant}
\author{}
\date{\today}

\begin{document}
\maketitle

\section{Notation and Data}

The legacy scripts are
\[
  \texttt{train\_network\_old.py},
  \qquad
  \texttt{train\_network\_old.sbatch}.
\]
The training set is made of synthetic pairs
\[
  \{(\data_i,\synth_i)\}_{i=1}^N,
\]
where each example contains two polarization channels on the same pixel grid
\(\grid\):
\[
  \data_i =
  \begin{pmatrix}
    d_{i,Q}\\ d_{i,U}
  \end{pmatrix}
  \in \R^{2\times 384\times 384},
  \qquad
  \synth_i =
  \begin{pmatrix}
    s'_{i,Q}\\ s'_{i,U}
  \end{pmatrix}
  \in \R^{2\times 384\times 384}.
\]
Here \(d\) is the contaminated/simulated observed map and \(s'\) is the clean
synthetic target map.  The target \(s'\) is drawn from the synthesis procedure
for the component-separated patch, then paired with a nuisance realization.

The legacy run uses all patches jointly:
\[
  p = 0,\ldots,191,
  \qquad
  \texttt{TRAIN\_FRACTION}=0.9,
  \qquad
  \texttt{N\_EPOCHS\_MEAN}=10,
  \qquad
  \texttt{N\_EPOCHS\_FULL}=90.
\]

\section{Normalization}

Before the network sees the data, both inputs and targets are standardized
channel-by-channel.  Let
\[
  m^{d}_{\chan},\; a^{d}_{\chan}
  \quad\text{and}\quad
  m^{s}_{\chan},\; a^{s}_{\chan},
  \qquad \chan\in\{Q,U\},
\]
be the empirical training-set mean and standard deviation for the data and
target channels.  The code computes these over all selected training files,
all augmentations, all syntheses, and all pixels.

The normalized tensors are
\[
  y_{i,\chan,\pix}
  =
  \frac{
    d_{i,\chan,\pix}-m^{d}_{\chan}
  }{
    a^{d}_{\chan}
  },
  \qquad
  x_{i,\chan,\pix}
  =
  \frac{
    s'_{i,\chan,\pix}-m^{s}_{\chan}
  }{
    a^{s}_{\chan}
  }.
\]
The network is trained in normalized units.  If it predicts a normalized mean
\(\widehat{\mean}_{\chan,\pix}\) and normalized standard deviation
\(\widehat{\std}_{\chan,\pix}\), the corresponding physical maps are
\[
  \mean_{\chan,\pix}
  =
  a^s_{\chan}\widehat{\mean}_{\chan,\pix}+m^s_{\chan},
  \qquad
  \std_{\chan,\pix}
  =
  a^s_{\chan}\widehat{\std}_{\chan,\pix}.
\]

\section{Legacy Network Parameterization}

The legacy code uses one shared U-Net
\[
  \net:\R^{2\times 384\times 384}
  \longrightarrow
  \R^{4\times 384\times 384}.
\]
Given normalized input \(y_i\), the four output channels are interpreted as
\[
  \net(y_i)
  =
  \left(
    \widehat{\mean}_{i,Q},
    \widehat{\mean}_{i,U},
    \widehat{\logvar}_{i,Q},
    \widehat{\logvar}_{i,U}
  \right).
\]
Equivalently, for each channel \(\chan\in\{Q,U\}\) and pixel
\(\pix\in\grid\),
\[
  \widehat{\mean}_{i,\chan,\pix}
  =
  [\net(y_i)]_{\chan,\pix},
  \qquad
  \widehat{\logvar}_{i,\chan,\pix}
  =
  [\net(y_i)]_{\chan+2,\pix}.
\]
The variance is parameterized through the log-variance:
\[
  \widehat{\std}_{i,\chan,\pix}^2
  =
  \exp\!\left(\widehat{\logvar}_{i,\chan,\pix}\right).
\]
The code clamps the log-variance before using it:
\[
  \widehat{\logvar}
  \leftarrow
  \operatorname{clip}(\widehat{\logvar},-12,12).
\]
This prevents numerically extreme variances.

\section{Legacy Stage 1: Mean Pretraining}

During the first stage, the same 4-output network is used, but only the first
two output channels enter the loss.  The log-variance channels are produced by
the network but ignored by the objective.

For a mini-batch \(\batch\), the loss is
\[
  \mathcal{L}_{\rm mean}(\theta)
  =
  \frac{1}{|\batch|\,2\,|\grid|}
  \sum_{i\in\batch}
  \sum_{\chan\in\{Q,U\}}
  \sum_{\pix\in\grid}
  \left(
    x_{i,\chan,\pix}
    -
    \widehat{\mean}_{\theta,\chan,\pix}(y_i)
  \right)^2 .
\]
This is an empirical conditional mean regression problem.  If the network class
and optimization were ideal, the minimizer would satisfy
\[
  \widehat{\mean}_{\theta}(y)
  \approx
  \E[x\mid y].
\]

\textbf{Important:} after this stage, the mean is not frozen in the legacy
script.  It is merely initialized by the MSE stage.

\section{Legacy Stage 2: Gaussian NLL}

In the second stage, the legacy script continues training the \emph{same}
network and the \emph{same} optimizer.  Now all four outputs enter the loss.
The code implements
\[
  \mathcal{L}_{\rm NLL}(\theta)
  =
  \frac{1}{|\batch|\,2\,|\grid|}
  \sum_{i\in\batch}
  \sum_{\chan\in\{Q,U\}}
  \sum_{\pix\in\grid}
  \frac{1}{2}
  \left[
    \widehat{\logvar}_{\theta,\chan,\pix}(y_i)
    +
    \frac{
      \left(
        x_{i,\chan,\pix}
        -
        \widehat{\mean}_{\theta,\chan,\pix}(y_i)
      \right)^2
    }{
      \exp\!\left(
        \widehat{\logvar}_{\theta,\chan,\pix}(y_i)
      \right)
    }
  \right],
\]
with the numerical modifications
\[
  \widehat{\logvar}
  \mapsto
  \operatorname{clip}(\widehat{\logvar},-12,12),
  \qquad
  \exp(\widehat{\logvar})
  \mapsto
  \max\{\exp(\widehat{\logvar}),10^{-8}\}.
\]

This is the negative log-likelihood of a diagonal Gaussian model in normalized
space:
\[
  x_{i,\chan,\pix}\mid y_i
  \sim
  \mathcal{N}
  \left(
    \widehat{\mean}_{\theta,\chan,\pix}(y_i),
    \exp\!\left[\widehat{\logvar}_{\theta,\chan,\pix}(y_i)\right]
  \right),
\]
where the covariance is diagonal over pixels and channels.  Thus, the std map
is a map of \emph{per-pixel marginal} posterior standard deviations, not a full
posterior covariance.

\section{What the Legacy Std Is Optimizing}

For a fixed input \(y_i\), channel \(\chan\), and pixel \(\pix\), define the
normalized residual
\[
  r_{i,\chan,\pix}(\theta)
  =
  x_{i,\chan,\pix}
  -
  \widehat{\mean}_{\theta,\chan,\pix}(y_i),
  \qquad
  v_{i,\chan,\pix}(\theta)
  =
  \exp\!\left[
    \widehat{\logvar}_{\theta,\chan,\pix}(y_i)
  \right].
\]
The local NLL contribution is
\[
  \ell_{\rm NLL}
  =
  \frac{1}{2}
  \left[
    \log v
    +
    \frac{r^2}{v}
  \right].
\]
If \(r\) is treated as fixed and we optimize only \(v\), then
\[
  \frac{\partial \ell_{\rm NLL}}{\partial v}
  =
  \frac{1}{2}
  \left(
    \frac{1}{v}
    -
    \frac{r^2}{v^2}
  \right).
\]
Setting this derivative to zero gives
\[
  v^\star = r^2.
\]
Equivalently,
\[
  \widehat{\logvar}^{\,\star}
  =
  \log r^2.
\]

This calculation explains the basic pressure on the std head: for fixed mean,
the Gaussian NLL wants the predicted variance to match the squared residual.
However, the legacy script does \emph{not} keep the mean fixed in stage 2.
Both
\[
  \widehat{\mean}_{\theta}
  \quad\text{and}\quad
  \widehat{\logvar}_{\theta}
\]
are functions of the same parameters \(\theta\), and the optimizer updates all
of them during the NLL stage.

\section{Why the Legacy Mean Can Change in Stage 2}

The gradient of the local NLL with respect to the mean is
\[
  \frac{\partial \ell_{\rm NLL}}{\partial \widehat{\mean}}
  =
  \frac{
    \widehat{\mean}-x
  }{
    v
  }.
\]
Thus, residuals in pixels with small predicted variance are weighted more
strongly.  The second stage therefore no longer solves ordinary MSE regression.
Instead, it solves a heteroscedastic weighted regression problem coupled to a
variance model:
\[
  \theta_{\rm full}
  =
  \arg\min_{\theta}
  \mathcal{L}_{\rm NLL}(\theta).
\]
The final legacy mean is therefore not necessarily the same as the stage-1 MSE
mean:
\[
  \widehat{\mean}_{\theta_{\rm full}}
  \neq
  \widehat{\mean}_{\theta_{\rm mean-stage}}
  \quad\text{in general}.
\]
This is the key point for interpreting the legacy std maps.  The legacy
variance is not trained against residuals of a frozen posterior mean.  It is
trained simultaneously with a mean that can move during the NLL stage.

\section{Interpretation of the Legacy Std Map}

The legacy std map represents
\[
  \widehat{\std}_{\theta,\chan,\pix}(d)
  =
  a^s_{\chan}
  \exp\!\left(
    \frac{1}{2}
    \widehat{\logvar}_{\theta,\chan,\pix}(y)
  \right),
\]
where \(y\) is the normalized version of \(d\).  Mathematically, this is a
diagonal approximation to
\[
  \operatorname{Std}\!\left[
    s_{\chan,\pix}\mid d
  \right].
\]
It is \emph{not}
\[
  \operatorname{Cov}(s_{\chan,\pix},s_{\chan',\pix'}\mid d),
\]
and it does not describe correlations across pixels or between \(Q\) and \(U\).

It is important to distinguish two different notions of ``separation''.  The
legacy NLL is channel-separated in the sense that it can be written as
\[
  \mathcal{L}_{\rm NLL}
  =
  \mathcal{L}_{Q}
  +
  \mathcal{L}_{U},
\]
with one Gaussian likelihood term for \(Q\) and one for \(U\).  However, the
\emph{function class} is not channel-separated.  The same U-Net trunk produces
the features used by both log-variance maps:
\[
  (d_Q,d_U)
  \xrightarrow{\text{shared trunk}}
  z_\theta(d_Q,d_U)
  \xrightarrow{\text{final }1\times1\text{ convolution}}
  \left(
    \log \std_Q^2,
    \log \std_U^2
  \right).
\]
At a given pixel \(\pix\), this is approximately a channel-specific linear
projection of the same shared feature vector:
\[
  \widehat{\logvar}_{Q,\pix}
  =
  a_Q^\top z_{\theta,\pix} + b_Q,
  \qquad
  \widehat{\logvar}_{U,\pix}
  =
  a_U^\top z_{\theta,\pix} + b_U.
\]
The vectors \(a_Q\) and \(a_U\) are different, but the spatial feature map
\(z_{\theta,\pix}\) is common.  Therefore, if the trunk learns a feature such
as
\[
  z_{1,\theta,\pix}
  \approx
  \text{local residual or noise amplitude at pixel }\pix,
\]
then both log-variance heads may primarily use this same feature:
\[
  \widehat{\logvar}_{Q,\pix}
  \approx
  \alpha_Q z_{1,\theta,\pix}+\beta_Q,
  \qquad
  \widehat{\logvar}_{U,\pix}
  \approx
  \alpha_U z_{1,\theta,\pix}+\beta_U.
\]
Even if \(\alpha_Q\neq \alpha_U\), the bright and dark regions of the two std
maps will be similar because both are functions of the same spatial template
\(z_{1,\theta,\pix}\).  The two maps may differ in amplitude while remaining
strongly correlated morphologically.

There is also a statistical reason why similarity can occur.  The variance is
driven by residual power.  For a fixed or slowly varying mean, the local NLL
pushes
\[
  \widehat{\std}_{Q,\pix}^2(d)
  \sim
  \left(s'_{Q,\pix}-\widehat{\mean}_{Q,\pix}(d)\right)^2,
  \qquad
  \widehat{\std}_{U,\pix}^2(d)
  \sim
  \left(s'_{U,\pix}-\widehat{\mean}_{U,\pix}(d)\right)^2.
\]
If the residual-power fields
\[
  \left(s'_Q-\widehat{\mean}_Q(d)\right)^2
  \quad\text{and}\quad
  \left(s'_U-\widehat{\mean}_U(d)\right)^2
\]
are large in the same regions, then similar std morphology is not necessarily a
network failure.  It may reflect shared patch geometry, common edge effects,
common synthesis artifacts, the same nuisance realization structure, or the
same non-stationarity across the patch.

Thus, strong correlation between
\[
  \widehat{\std}_Q(d)
  \quad\text{and}\quad
  \widehat{\std}_U(d)
\]
can have two sources:
\[
  \text{architectural coupling through }z_\theta(d),
  \qquad
  \text{and/or}
  \qquad
  \text{similar residual-power targets}.
\]
Separating the std trunks would address only the architectural coupling.  If
the residual-power targets are themselves similar, independent trunks could
still learn similar std maps.

\section{Latest Legacy-Style Variant: Separated Std Trunks}

The latest variant keeps the legacy probabilistic objective, but changes the
network parameterization used for the standard deviations.  The change is not a
change in the likelihood.  It is a change in the function class.

In the previous single-trunk legacy architecture, one U-Net produced all four
posterior maps:
\[
  F_\theta(y)
  =
  \left(
    \widehat{\mean}_{Q}(y),
    \widehat{\mean}_{U}(y),
    \widehat{\logvar}_{Q}(y),
    \widehat{\logvar}_{U}(y)
  \right).
\]
Thus the posterior means and posterior log-variances were all functions of the
same learned representation:
\[
  y
  \xrightarrow{\text{shared U-Net trunk}}
  z_\theta(y)
  \xrightarrow{\text{output head}}
  \left(
    \widehat{\mean}_{Q},
    \widehat{\mean}_{U},
    \widehat{\logvar}_{Q},
    \widehat{\logvar}_{U}
  \right).
\]

The latest legacy-style variant instead uses one mean trunk and two independent
std trunks:
\[
  M_\theta(y)
  =
  \left(
    \widehat{\mean}_{Q}(y),
    \widehat{\mean}_{U}(y)
  \right),
\]
\[
  V_{\theta_Q}(y)
  =
  \widehat{\logvar}_{Q}(y),
  \qquad
  V_{\theta_U}(y)
  =
  \widehat{\logvar}_{U}(y).
\]
Equivalently, the full model is
\[
  F_{\theta,\theta_Q,\theta_U}(y)
  =
  \left(
    M_\theta(y),
    V_{\theta_Q}(y),
    V_{\theta_U}(y)
  \right),
\]
or, componentwise,
\[
  F_{\theta,\theta_Q,\theta_U}(y)
  =
  \left(
    \widehat{\mean}_{Q}(y),
    \widehat{\mean}_{U}(y),
    \widehat{\logvar}_{Q}(y),
    \widehat{\logvar}_{U}(y)
  \right).
\]

The two-stage training structure remains legacy-like.  First, the mean trunk is
trained by the MSE objective
\[
  \mathcal{L}_{\rm mean}(\theta)
  =
  \frac{1}{|\batch|\,2\,|\grid|}
  \sum_{i\in\batch}
  \sum_{\chan\in\{Q,U\}}
  \sum_{\pix\in\grid}
  \left[
    x_{i,\chan,\pix}
    -
    \widehat{\mean}_{\theta,\chan,\pix}(y_i)
  \right]^2.
\]
Then the full model is trained with the same diagonal Gaussian NLL as before:
\[
  \mathcal{L}_{\rm NLL}(\theta,\theta_Q,\theta_U)
  =
  \frac{1}{|\batch|\,2\,|\grid|}
  \sum_{i\in\batch}
  \sum_{\chan\in\{Q,U\}}
  \sum_{\pix\in\grid}
  \frac{1}{2}
  \left[
    \widehat{\logvar}_{\chan,\pix}(y_i)
    +
    \frac{
      \left(
        x_{i,\chan,\pix}
        -
        \widehat{\mean}_{\chan,\pix}(y_i)
      \right)^2
    }{
      \exp\!\left(\widehat{\logvar}_{\chan,\pix}(y_i)\right)
    }
  \right].
\]
Here
\[
  \widehat{\logvar}_{Q}=V_{\theta_Q}(y),
  \qquad
  \widehat{\logvar}_{U}=V_{\theta_U}(y),
\]
while
\[
  (\widehat{\mean}_Q,\widehat{\mean}_U)=M_\theta(y).
\]

Therefore the difference between the previous legacy architecture and the
latest variant is
\[
  \text{previous:}\qquad
  (\widehat{\mean}_Q,\widehat{\mean}_U,
   \widehat{\logvar}_Q,\widehat{\logvar}_U)
  =
  F_\theta(y),
\]
versus
\[
  \text{latest:}\qquad
  (\widehat{\mean}_Q,\widehat{\mean}_U)=M_\theta(y),
  \qquad
  \widehat{\logvar}_Q=V_{\theta_Q}(y),
  \qquad
  \widehat{\logvar}_U=V_{\theta_U}(y).
\]

The likelihood is still factorized over channels and pixels:
\[
  p(x\mid y)
  \approx
  \prod_{\pix\in\grid}
  p(x_{Q,\pix}\mid y)\,
  p(x_{U,\pix}\mid y).
\]
Thus the std maps remain maps of per-pixel marginal posterior standard
deviations:
\[
  \widehat{\std}_{Q,\pix}(d)
  \approx
  \operatorname{Std}[s_{Q,\pix}\mid d],
  \qquad
  \widehat{\std}_{U,\pix}(d)
  \approx
  \operatorname{Std}[s_{U,\pix}\mid d].
\]
The model still does not learn the full posterior covariance
\[
  \operatorname{Cov}(s_{\chan,\pix},s_{\chan',\pix'}\mid d).
\]

The purpose of the separated std trunks is to remove the architectural pressure
that forces \(\widehat{\logvar}_Q\) and \(\widehat{\logvar}_U\) to be read out
from the same trunk representation \(z_\theta(y)\).  The Q std map now has its
own parameters \(\theta_Q\), and the U std map has its own parameters
\(\theta_U\).  This gives the two uncertainty maps more independent degrees of
freedom.  It does not, however, force them to be different: if the Q and U
residual-power fields are statistically similar, the independently trained
trunks may still learn morphologically similar std maps.

\section{Current Version: Fixed Legacy Mean and New Std Network}

The current version introduces a different factorization.  There are two
networks:
\[
  g_\phi(y)
  =
  \widehat{\mean}_{\phi}(y)
  \in \R^{2\times384\times384},
  \qquad
  h_\psi(y)
  =
  \widehat{\logvar}_{\psi}(y)
  \in \R^{2\times384\times384}.
\]
In the specific current run using the legacy mean,
\[
  g_\phi
\]
is initialized from
\[
  \texttt{version\_2/models/moment\_network\_joint\_final.pth}.
\]
The old legacy checkpoint had four outputs:
\[
  (\widehat{\mean}_Q,\widehat{\mean}_U,
    \widehat{\logvar}_Q,\widehat{\logvar}_U).
\]
The current script copies only the first two output channels:
\[
  g_\phi(y)
  =
  \left[
    f_{\theta_{\rm legacy}}(y)
  \right]_{0:2}.
\]
The old legacy log-variance channels are discarded.

The sbatch sets
\[
  \texttt{N\_EPOCHS\_MEAN}=0,
  \qquad
  \texttt{N\_EPOCHS\_FULL}=90,
  \qquad
  \texttt{USE\_INTENSITY}=0.
\]
Therefore no mean training is run in the current job.  During std training,
the code evaluates
\[
  \widehat{\mean}_{\phi}(y_i)
\]
inside a \(\texttt{torch.no\_grad()}\) block, and the optimizer is applied only
to the parameters \(\psi\) of \(h_\psi\).  Hence
\[
  \phi \quad\text{is fixed during the std stage}.
\]

\section{Current Std Objective}

The current std stage does not use the Gaussian NLL directly.  Instead, with
the fixed mean \(g_\phi\), it defines a residual target
\[
  R_{i,\chan,\pix}
  =
  \left(
    x_{i,\chan,\pix}
    -
    \widehat{\mean}_{\phi,\chan,\pix}(y_i)
  \right)^2,
\]
then forms
\[
  T_{i,\chan,\pix}
  =
  \log\left(
    \max\{R_{i,\chan,\pix},\epsilon\}
  \right),
  \qquad
  \epsilon=\texttt{RESIDUAL\_FLOOR}=10^{-8}.
\]
The variance network predicts
\[
  \widehat{\logvar}_{\psi,\chan,\pix}(y_i)
  =
  h_{\psi,\chan,\pix}(y_i),
\]
clamped to \([-12,12]\), and the loss is
\[
  \mathcal{L}_{\rm std}(\psi)
  =
  \frac{1}{|\batch|\,2\,|\grid|}
  \sum_{i\in\batch}
  \sum_{\chan\in\{Q,U\}}
  \sum_{\pix\in\grid}
  \left[
    \operatorname{clip}
    \left(
      h_{\psi,\chan,\pix}(y_i),-12,12
    \right)
    -
    \operatorname{clip}
    \left(
      T_{i,\chan,\pix},-12,12
    \right)
  \right]^2.
\]
Thus, the current run is a direct regression onto the log squared residuals of
the fixed mean model:
\[
  h_\psi(y)
  \approx
  \log\left[
    \left(x-g_\phi(y)\right)^2
  \right].
\]

\section{Main Difference}

The mathematical distinction is:
\[
\begin{array}{lll}
\toprule
 & \text{Legacy} & \text{Current legacy-mean/std run} \\
\midrule
\text{Mean model} &
\text{same network as std} &
\text{loaded from legacy checkpoint} \\
\text{Mean during std training} &
\text{updated by NLL} &
\text{fixed} \\
\text{Std model} &
\text{channels 2:4 of same network} &
\text{separate U-Net} \\
\text{Std objective} &
\text{Gaussian NLL} &
\text{MSE on log squared residuals} \\
\text{Input channels} &
(d_Q,d_U) &
(d_Q,d_U) \\
\bottomrule
\end{array}
\]

In compact form, the legacy full stage solves
\[
  \min_{\theta}
  \sum_{i,\chan,\pix}
  \frac{1}{2}
  \left[
    \widehat{\logvar}_{\theta}
    +
    \frac{
      (x-\widehat{\mean}_{\theta})^2
    }{
      \exp(\widehat{\logvar}_{\theta})
    }
  \right],
\]
whereas the current legacy-mean/std run solves
\[
  \min_{\psi}
  \sum_{i,\chan,\pix}
  \left[
    h_\psi(y_i)
    -
    \log
    \left(
      (x_i-g_\phi(y_i))^2
    \right)
  \right]^2,
  \qquad
  \phi \text{ fixed}.
\]

\section{Practical Consequence}

If the goal is to keep the legacy mean exactly as the reference mean and only
modify the uncertainty model, the current setup is closer to that goal:
\[
  \widehat{\mean}_{\rm current}(d)
  =
  \widehat{\mean}_{\rm legacy}(d),
  \qquad
  \text{only } \widehat{\std}(d) \text{ is newly trained}.
\]
In the checked local outputs, the normalization statistics of
\(\texttt{version\_2}\) and
\(\texttt{version\_2\_qu\_legacymean\_std}\) match exactly, so the copied
legacy mean is evaluated in the same normalized coordinates.

However, the std objective is no longer the same as the legacy NLL.  It is a
direct regression to log residual power.  This makes the interpretation clearer
for the fixed-mean experiment, but it is not mathematically identical to the
joint maximum-likelihood training used by the legacy script.

\end{document}
